
\documentclass[12pt, a4paper]{article}

\setlength{\headheight}{14.32365pt}
\addtolength{\topmargin}{-0.00365pt}

\usepackage[a4paper, top=2.5cm, bottom=2.5cm, left=2.5cm, right=2.5cm]{geometry}
\usepackage[utf8]{inputenc}
\usepackage[T1]{fontenc}
\usepackage[spanish, es-tabla]{babel}
\usepackage{helvet} % Tipografía sans-serif limpia estilo Helvetica/Arial
\renewcommand{\familydefault}{\sfdefault} % Configura el documento en sans-serif


\babelprovide[import, onchar=ids fonts]{spanish}
\babelprovide[import, onchar=ids fonts]{english}


\usepackage{enumitem}
\setlist[itemize]{label=-}

\usepackage{xcolor}
% Azul Oficial ITBA (#002B49)
\definecolor{itbablue}{RGB}{0, 43, 73}
\definecolor{itbagris}{RGB}{130, 138, 145}

\usepackage{graphicx}
\usepackage{tikz}
\usetikzlibrary{calc, positioning, arrows.meta, decorations.pathreplacing}
\usepackage{pgfplots}
\pgfplotsset{compat=1.18}
\usepackage{float}
\usepackage{titlesec}
\usepackage{fancyhdr}
\usepackage{setspace}

% Configuración de estilo de títulos de sección
\titleformat{\section}
  {\normalfont\Large\bfseries\color{itbablue}}{\thesection.}{0.5em}{}
\titleformat{\subsection}
  {\normalfont\large\bfseries\color{itbablue}}{\thesubsection.}{0.5em}{}

% Configuración de encabezados y pies de página para páginas interiores
\pagestyle{fancy}
\fancyhf{}
\rhead{\small\textcolor{gray}{ITBA -- Proyecto de Tesis}}
\lhead{\small\textcolor{gray}{\leftmark}}
\rfoot{\thepage}
\renewcommand{\headrulewidth}{0.4pt}

\usepackage{hyperref}
\hypersetup{
    colorlinks=true,
    linkcolor=itbablue,
    citecolor=itbablue,
    urlcolor=itbablue,
    pdfauthor={Moccagatta, Nicolás},
    pdftitle={Efecto causal de las migraciones de tecnología frontend sobre la experiencia de usuario real -- ITBA}
}

\begin{document}

\begin{titlepage}
    % Dibujo del marco perimetral idéntico al PDF original
    \begin{tikzpicture}[remember picture, overlay]
        \draw[line width=1.2pt, color=itbablue]
            ($(current page.north west) + (1.2cm,-1.2cm)$)
            rectangle
            ($(current page.south east) + (-1.2cm,1.2cm)$);
        \draw[line width=0.5pt, color=itbablue]
            ($(current page.north west) + (1.3cm,-1.3cm)$)
            rectangle
            ($(current page.south east) + (-1.3cm,1.3cm)$);
    \end{tikzpicture}

    \vspace*{-0.2cm}

    % Encabezado superior: Nombre de la institución a la izquierda y Logo ITBA a la derecha
    \noindent
    \begin{minipage}[c]{0.6\textwidth}
        {\bfseries\large INSTITUTO TECNOLÓGICO\\DE BUENOS AIRES}
    \end{minipage}%
    \begin{minipage}[c]{0.4\textwidth}
        \raggedleft
        \IfFileExists{logo-itba.png}{%
            \includegraphics[width=4.2cm]{logo-itba.png}%
        }{%
            \framebox{\parbox{3.8cm}{\centering \textbf{ITBA}\\[2pt]\small Logo Institucional}}%
        }
    \end{minipage}

    \vfill

    % Título del proyecto centrado
    \begin{center}
        {\fontsize{16pt}{24pt}\selectfont \textbf{EFECTO CAUSAL DE LAS MIGRACIONES DE TECNOLOGÍA FRONTEND SOBRE LA EXPERIENCIA DE USUARIO REAL: UN ESTUDIO CUASI-EXPERIMENTAL A GRAN ESCALA SOBRE CORE WEB VITALS}}
    \end{center}

    \vfill

    % Información del trabajo y de la materia
    \noindent
    \begin{minipage}{\textwidth}
        \setstretch{1.3}
        \large
        \textbf{ALUMNO:} Moccagatta, Nicolás (Leg. N.\textsuperscript{o} 106396) \\[0.3cm]
        \textbf{DIRECTOR DE TESIS:} Salgado Corrado, Ariel Olaf \\[0.3cm]
        \textbf{CO-DIRECTOR DE TESIS:} Sarraute, Carlos \\[0.3cm]
        \textbf{TALLER:} Taller de Tesis 1 -- Maestría en Ciencia de Datos \\[0.3cm]
        \textbf{PROFESORES:} Salgado Corrado, Ariel Olaf; Sarraute, Carlos
    \end{minipage}

    \vfill
    \vspace{0.8cm}

    % Pie de la carátula centrado
    \begin{center}
        {\large \textbf{BUENOS AIRES}} \\[0.2cm]
        {\large \textbf{SEGUNDO CUATRIMESTRE, 2026}}
    \end{center}
    \vspace*{0.3cm}
\end{titlepage}

\newpage
\tableofcontents
\newpage




\section{Introducción}

\subsection{Contexto del proyecto}

Cuando un usuario abre un sitio web en su celular o su computadora, la velocidad con la que ese sitio responde condiciona directamente si se queda, si compra, si vuelve. Google intentó estandarizar esa noción difusa de ``velocidad percibida'' en un conjunto de tres métricas llamadas Core Web Vitals: el Largest Contentful Paint (LCP), que mide cuánto tarda en aparecer el elemento visual más grande de la pantalla; el Interaction to Next Paint (INP), que mide qué tan rápido responde la página a una interacción del usuario; y el Cumulative Layout Shift (CLS), que mide cuánto se mueve el contenido de forma inesperada mientras la página carga. Estas tres métricas se recolectan de navegadores Chrome reales, de forma agregada y anónima, y se publican mensualmente en un dataset público llamado Chrome User Experience Report (CrUX). A diferencia de un benchmark de laboratorio, CrUX refleja la experiencia efectiva de miles de millones de visitas reales.

En paralelo, existe otro proyecto público llamado HTTP Archive, que rastrea mensualmente millones de sitios web y, entre otras cosas, detecta qué tecnologías usa cada uno: si corre sobre WordPress o Shopify, si el frontend está armado con jQuery, React, Vue o algún framework más reciente como Next.js o Astro. Esta detección se hace con una herramienta llamada Wappalyzer, que identifica tecnologías a partir de patrones en el código fuente, las cabeceras HTTP y otros indicios técnicos.

HTTP Archive y CrUX están relacionados pero son dos fuentes distintas, mantenidas por equipos distintos. Desde 2018, HTTP Archive usa la lista de orígenes de CrUX para decidir qué sitios rastrear, y cada página que rastrea incluye además una copia puntual de los valores de CrUX vigentes ese mes; el panel mensual completo que usa este trabajo se toma directamente de la publicación propia de CrUX, no de esa copia embebida. HTTP Archive se usa exclusivamente como fuente de la tecnología detectada en cada sitio, que es el dato que CrUX no tiene.

La industria del desarrollo web toma constantemente decisiones de migración entre estas tecnologías (dejar jQuery por React, pasar de WordPress a una arquitectura desacoplada tipo Jamstack) y suele justificarlas, al menos en parte, con la promesa de una mejora en performance. Sin embargo, esa promesa rara vez se contrasta con evidencia rigurosa: lo habitual es apoyarse en casos de éxito puntuales de una empresa, en benchmarks de laboratorio sobre aplicaciones de prueba, o en comparaciones de sitios distintos en un mismo momento, que confunden el efecto de la tecnología con otros factores (presupuesto del equipo, tipo de negocio, madurez del sitio). El cruce de HTTP Archive y CrUX abre una posibilidad distinta: seguir a los mismos sitios en el tiempo y observar qué pasa con su experiencia de usuario real cuando efectivamente cambian de tecnología.

Esta pregunta viene madurando desde antes del inicio del taller. En un trabajo de la cursada de la maestría se cruzó, por curiosidad personal, el historial de escucha musical del propio alumno con su actividad de commits en GitHub, para ver si existía alguna relación entre ambos. El análisis se limitó deliberadamente a explorar co-ocurrencia y no a probar causalidad, porque el diseño disponible en ese momento (dos series de tiempo superpuestas) no alcanzaba para más. Ese límite, declarado en su momento, es el punto de partida de este trabajo: acá se dispone de un evento de tratamiento bien definido (la migración tecnológica), un grupo de sitios que no lo recibe, y una batería de métodos pensados específicamente para separar una asociación de un efecto.

\subsection{Justificación del estudio}

La justificación de este trabajo parte de tres limitaciones del conocimiento disponible hoy sobre la relación entre tecnología frontend y performance web.

En primer lugar, la evidencia que produce la propia industria (el Core Web Vitals Technology Report de HTTP Archive \cite{httparchive2025cwvtech}, el capítulo de Performance del Web Almanac \cite{httparchive2025almanac}) es transversal: compara, en un corte del tiempo, qué porcentaje de sitios construidos con cada tecnología aprueba los Core Web Vitals. Ese tipo de comparación no controla por el hecho de que los sitios que usan una tecnología u otra pueden ser sistemáticamente distintos entre sí desde antes de adoptarla.

En segundo lugar, la literatura académica sobre frameworks frontend es, en el mejor de los casos, experimental a pequeña escala sobre aplicaciones de prueba armadas específicamente para el estudio \cite{tong2023comparative}, y en el resto de los casos es una síntesis narrativa de literatura y casos reportados por la industria, sin datos propios \cite{tulsyan2024impact}. Ninguna de las dos aporta lo que hace falta para responder esta pregunta: qué pasa cuando un sitio real, con su tráfico real, cambia de tecnología.

En tercer lugar, ningún trabajo identificado combina las tres condiciones necesarias para responder esta pregunta con solidez: seguimiento longitudinal de los mismos sitios en el tiempo, una medida de experiencia de usuario real (no de laboratorio) y un diseño que permita una lectura causal, en vez de puramente correlacional. La revisión sistemática más reciente sobre predicción de performance web, que relevó 6.657 estudios entre 2010 y 2024 \cite{ghattas2025review}, confirma este vacío: encuentra que el campo se concentra en modelos que predicen un puntaje de performance a partir de métricas técnicas, mayormente en sitios de gobierno electrónico y educación, y no identifica estudios que estimen el efecto de decisiones tecnológicas concretas sobre la experiencia de usuarios reales.

El aporte de este trabajo se ubica exactamente en ese cruce: construir un panel longitudinal de sitios a partir de datos públicos y aplicar técnicas de inferencia causal, ya maduras en economía aplicada pero prácticamente no usadas en el estudio de la performance web, para estimar si las migraciones de tecnología realmente producen la mejora que se les atribuye, y en qué casos. En el plano científico, esto llena un vacío metodológico concreto. En el plano práctico, aporta evidencia a un tipo de decisión que hoy la industria toma con muy poco sustento empírico.

Vale mencionar, además, que la elección de este tema responde también a una razón práctica aprendida en trabajos previos de la maestría: los datos de HTTP Archive y de CrUX son completamente públicos y accesibles vía BigQuery, sin necesidad de tramitar autorizaciones institucionales ni de trabajar con datos de personas, a diferencia de otras fuentes (redes sociales) exploradas anteriormente y descartadas por restricciones de acceso.

Buena parte de la ingeniería de datos que requiere este trabajo ya fue practicada, aunque sobre otro dominio, en materias previas de la maestría: la extracción de tablas particionadas por fecha desde un dataset público de BigQuery hacia un esquema dimensional propio, el versionado de atributos que cambian en el tiempo (dimensiones de cambio lento, o SCD) y la corrección de discontinuidades de medición entre sucesivas mediciones de un panel. Lo que no está cubierto por ninguna materia cursada, y constituye el aporte metodológico central de esta tesis, es el armazón de inferencia causal propiamente dicho: el diseño de diferencias en diferencias, el tratamiento de la adopción escalonada y la estimación de heterogeneidad de efectos que se detallan en la Sección \ref{sec:solucion}.

\subsection{Alcances del trabajo}

El estudio se circunscribe a un conjunto acotado de pares de migración tecnológica (por ejemplo, jQuery a React, o el abandono de WordPress), elegidos por tener volumen suficiente de casos y una detección estable en Wappalyzer a lo largo del tiempo. La ventana temporal utilizable es aproximadamente 2018-2026, período en el que el corpus de HTTP Archive ya sigue la lista de orígenes de CrUX de forma consistente. La cohorte de sitios elegibles se define de forma point-in-time (según su estado antes del evento de migración), para no introducir sesgo de supervivencia.

El diseño es observacional y cuasi-experimental, no un experimento aleatorizado: se estima el efecto promedio sobre los sitios tratados (ATT, average treatment effect on the treated) y su heterogeneidad, sin pretensión de generalizar a cualquier migración imaginable ni de establecer causalidad en el sentido experimental estricto.

\section{Estado de la cuestión}

La literatura relevante se ordena en dos bloques. El primero son los antecedentes de dominio, es decir, lo que ya se investigó sobre tecnología web y performance, y se abre en tres líneas: la evidencia industrial que cruza tecnología con Core Web Vitals, la evidencia experimental de laboratorio sobre frameworks frontend y el uso descriptivo de CrUX a gran escala. El segundo bloque son los antecedentes metodológicos: la literatura de inferencia causal sobre datos de panel con tratamiento escalonado, ajena al dominio de performance web pero central para el diseño de este trabajo.

\subsection{Evidencia industrial y de laboratorio}

Una primera línea es la evidencia de tipo industrial sobre tecnología y Core Web Vitals: el Core Web Vitals Technology Report \cite{httparchive2025cwvtech} y el Web Almanac \cite{httparchive2025almanac}, ambos producidos por HTTP Archive, cruzan la tecnología detectada en cada sitio con su desempeño en CrUX. Es la fuente de datos más parecida a la que usa este trabajo, pero su análisis es descriptivo y transversal: compara sitios distintos en un mismo mes, sin seguirlos en el tiempo ni controlar por diferencias previas entre ellos.

Una segunda línea es la literatura sobre frameworks frontend y su desempeño. Tong, Jikson y Gunawan \cite{tong2023comparative} comparan React, Vue y Next.js sobre aplicaciones de referencia construidas específicamente para el estudio, midiendo tiempos de respuesta y el costo de operaciones puntuales sobre el DOM; es evidencia experimental limpia, pero de pequeña escala y sin usuarios reales de por medio. Tulsyan et al. \cite{tulsyan2024impact}, en cambio, no corren un experimento propio: sintetizan literatura existente y casos reportados por la industria sobre el impacto de los frameworks de JavaScript en la performance y la experiencia de usuario, con un aporte panorámico más que cuantitativo. Ninguno de los dos trabajos sigue sitios reales en producción a lo largo del tiempo ni usa datos de usuarios.

\subsection{Uso descriptivo de CrUX y vacío metodológico}

Una tercera línea usa CrUX a gran escala con fines descriptivos. Sengupta et al. \cite{sengupta2023crux} analizan la experiencia de navegación real de los usuarios de CrUX según el dispositivo y el país, mostrando el potencial del dataset para investigación a escala. El trabajo, sin embargo, no estudia tecnología ni migraciones: caracteriza la experiencia de navegación en general.

A esto se suma la revisión sistemática de Ghattas, Odeh y Mora \cite{ghattas2025review}, que sintetiza treinta estudios de alta calidad sobre predicción de performance web publicados entre 2010 y 2024, y concluye que el campo se concentra en modelos de predicción de un puntaje de performance, mayormente sobre sitios de gobierno y educación, con escasa adopción de aprendizaje automático y sin estudios de efectos causales de decisiones de arquitectura.

\subsection{Inferencia causal para datos de panel con tratamiento escalonado}

Una cuarta línea, ajena al dominio de performance web pero central para la metodología de este trabajo, es la literatura de inferencia causal sobre datos de panel con tratamiento escalonado (staggered treatment), es decir, cuando distintas unidades reciben un tratamiento (en este caso, la migración) en distintos momentos del tiempo. El diseño clásico de diferencias en diferencias (DiD) se remonta al trabajo de Card y Krueger \cite{cardkrueger1994}, que comparó el empleo en locales de comida rápida de dos estados con distinta política de salario mínimo, antes y después del cambio. Cuando el tratamiento no ocurre en un único momento sino escalonado en el tiempo, como es el caso de las migraciones tecnológicas, el estimador clásico de efectos fijos bidireccionales (TWFE) puede arrojar estimaciones sesgadas, porque termina comparando implícitamente unidades que ya recibieron el tratamiento contra otras que lo recibieron antes, en vez de contra un grupo de control no tratado; este problema fue formalizado por Goodman-Bacon \cite{goodmanbacon2021}. Callaway y Sant'Anna \cite{callawaysantanna2021} proponen un estimador robusto a ese sesgo, que identifica el efecto del tratamiento separadamente para cada cohorte de adopción y cada período, y lo agrega de forma consistente. Para estudiar si el efecto de una migración varía según las características del sitio, existe además una familia de métodos de aprendizaje automático causal, entre los que se destacan los bosques aleatorios causales (causal forests) de Wager y Athey \cite{wagerathey2018}, que permiten estimar el efecto de tratamiento condicional a covariables (CATE) sin imponer una forma funcional arbitraria.

En síntesis, el vacío que aborda este trabajo se ubica en la intersección de estas líneas: existe evidencia industrial transversal sobre tecnología y Core Web Vitals, existe evidencia experimental de laboratorio sobre frameworks, existe uso descriptivo de CrUX a escala y existe un conjunto de métodos causales maduro en econometría aplicada, pero nadie combinó las cuatro cosas para estimar el efecto de las decisiones de arquitectura de software sobre la experiencia real de los usuarios.

\section{Definición del problema}

\subsection{Problema general}

El problema que aborda este trabajo es estimar el efecto causal de las migraciones de tecnología frontend, detectadas en el panel mensual de HTTP Archive, sobre las métricas de experiencia de usuario real publicadas en CrUX, tratando cada migración como un cuasi-experimento y aplicando un diseño de diferencias en diferencias robusto a tratamiento escalonado, con un grupo de control de sitios comparables que no migraron, utilizando exclusivamente datos públicos.

Se distingue de los antecedentes más cercanos en tres aspectos. Respecto del Core Web Vitals Technology Report y el Web Almanac, la diferencia es de diseño: en vez de comparar sitios distintos en un corte del tiempo, este trabajo sigue a los mismos sitios antes y después de un evento de migración. Respecto de los estudios de laboratorio, la diferencia es de escala y de validez externa: en vez de aplicaciones de prueba, se trabaja con el universo de sitios reales que cubre CrUX. Respecto del uso descriptivo de CrUX en la literatura, la diferencia es de objetivo: en vez de caracterizar la experiencia de navegación en general, se aísla el efecto de un tipo específico de decisión, la migración tecnológica.

\subsection{Preguntas de investigación}

¿Las migraciones desde tecnologías frontend con mayor carga de JavaScript bloqueante hacia frameworks modernos producen una mejora causal y estadísticamente significativa en la experiencia de usuario real (LCP, INP), respecto de la trayectoria que hubieran seguido los mismos sitios de no migrar? ¿Ese efecto es homogéneo entre distintos tipos de migración y niveles de tráfico, o varía sistemáticamente según las características del sitio? ¿En qué medida un estimador clásico de efectos fijos bidireccionales, ampliamente usado en la práctica, sobreestima ese efecto frente a un estimador robusto al sesgo de tratamiento escalonado?

\section{Descripción del set de datos}
\label{sec:dataset}

\subsection{Origen y contexto del dataset}

Los datos que sostienen este trabajo son en su totalidad públicos y se acceden vía BigQuery: la tabla de crawls de HTTP Archive (\texttt{httparchive.crawl.pages}), que incluye para cada origen y cada mes la tecnología detectada por Wappalyzer, y la tabla materializada de CrUX (\texttt{chrome-ux-report.materialized.metrics\_summary}), que incluye para cada origen y cada mes el percentil 75 de LCP, INP y CLS. No se trabaja en ningún momento con datos personales: la unidad de análisis es el origen web (el dominio), no el usuario individual.

El uso de ambas fuentes con fines de investigación tiene precedente en la literatura, aunque no en la combinación causal que propone este trabajo. Ruth et al. \cite{ruth2022toppling}, en un estudio publicado en la Internet Measurement Conference de ACM, evaluaron la exactitud de distintos rankings de sitios web y concluyeron que CrUX es, de las fuentes comparadas, la que mejor refleja la popularidad real de un sitio \cite{chromedev2025cruxmethodology}, lo que respalda su uso como dataset de investigación. Beugin et al. \cite{beugin2026lessons} combinaron crawls históricos de HTTP Archive con datos de CrUX para estudiar, de forma longitudinal a lo largo de siete años, la adopción y el abandono de un conjunto de APIs web, en un diseño de datos estructuralmente similar al que se usa en este trabajo, aunque con un objetivo distinto (caracterizar la adopción de una tecnología, no estimar su efecto causal sobre una métrica de resultado).

\subsection{Descripción técnica}

Un relevamiento inicial contra ambas fuentes (BigQuery, agosto de 2026) confirma la viabilidad del panel, tomando julio de 2026 como mes de referencia común para ambas fuentes. El outcome de CrUX está disponible de forma mensual entre 2017-10 y 2026-07; en julio de 2026 cubre \textbf{18.059.068} origins. HTTP Archive, filtrando únicamente páginas raíz para alinear la granularidad con CrUX, releva \textbf{16.268.243} orígenes ese mismo mes. La cobertura y distribución de calidad de las tres métricas en julio de 2026 son: LCP, 97,3\,\% de cobertura, 70,8\,\% ``Good'' ($\leq$2,5\,s), 19,8\,\% ``Needs Improvement'' (2,5-4\,s), 9,4\,\% ``Poor'' ($>$4\,s); INP, 70,1\,\% de cobertura (menor porque requiere interacciones del usuario), 88,4\,\% Good ($\leq$200\,ms), 9,5\,\% Needs Improvement (200-500\,ms), 2,0\,\% Poor ($>$500\,ms); CLS, 99,4\,\% de cobertura, 87,2\,\% Good ($\leq$0,10), 6,5\,\% Needs Improvement (0,10-0,25), 6,3\,\% Poor ($>$0,25). La Figura \ref{fig:semaforo} resume esa distribución.

\begin{figure}[H]
\centering
\begin{tikzpicture}
\begin{axis}[
    xbar stacked,
    width=0.78\textwidth, height=4.4cm,
    xmin=0, xmax=100,
    xlabel={\% de orígenes que tienen dato en la métrica},
    xlabel style={font=\small},
    ytick=data,
    symbolic y coords={CLS,INP,LCP},
    y tick label style={font=\small},
    x tick label style={font=\scriptsize},
    bar width=14pt,
    axis lines=left,
    legend style={at={(0.5,-0.45)}, anchor=north, legend columns=3,
                  font=\scriptsize, draw=gray!40},
]
\addplot[fill=itbablue, draw=none] coordinates {(70.8,LCP) (88.4,INP) (87.2,CLS)};
\addplot[fill=itbablue!35, draw=none] coordinates {(19.8,LCP) (9.5,INP) (6.5,CLS)};
\addplot[fill=itbagris, draw=none] coordinates {(9.4,LCP) (2.0,INP) (6.3,CLS)};
\legend{Good, Needs Improvement, Poor}
\end{axis}
\end{tikzpicture}
\caption{Distribución de calidad de las tres métricas de Core Web Vitals en CrUX, julio de 2026, según los umbrales oficiales de Google. La base de cada barra son los orígenes que tienen dato en esa métrica: 97,3\,\% del total para LCP, 70,1\,\% para INP y 99,4\,\% para CLS.}
\label{fig:semaforo}
\end{figure}

El volumen de eventos de migración entre dos cortes separados por un año (junio de 2025 y junio de 2026), contando únicamente orígenes de páginas raíz presentes en el crawl móvil de ambos meses, es de 62.457 migraciones de jQuery a React, 65.407 abandonos de WordPress y 9.059 migraciones de React a Next.js. Como se trata de los eventos visibles comparando solo dos meses, el conteo sobre el panel mensual completo será necesariamente mayor. Aun así, ya alcanza con holgura para un diseño de diferencias en diferencias con cohortes de adopción escalonada, que necesita del orden de cientos de unidades tratadas por cohorte.

La Tabla \ref{tab:variables} lista las variables operativas del panel y el papel que cumple cada una en el diseño. La unidad de observación es el par origen-mes; la tecnología detectada se versiona como una dimensión de cambio lento (SCD tipo 2, con rango de vigencia desde/hasta), de modo que el mes de inicio de vigencia de un valor nuevo es la fecha del evento de tratamiento.

\begin{table}[htbp]
\centering
\small
\caption{Variables operativas del panel.}
\label{tab:variables}
\begin{tabular}{|p{3.3cm}|p{3.2cm}|p{2.6cm}|p{4.1cm}|}
\hline
\textbf{Variable} & \textbf{Fuente} & \textbf{Granularidad} & \textbf{Rol en el diseño} \\ \hline
Tecnología detectada (Wappalyzer) & HTTP Archive, \texttt{crawl.pages} & Origen, mes y cliente & Define el tratamiento: su cambio sostenido es el evento de migración \\ \hline
p75 de LCP & CrUX, \texttt{metrics\_summary} & Origen y mes & Variable de resultado principal (H1) \\ \hline
p75 de INP & CrUX, \texttt{metrics\_summary} & Origen y mes & Variable de resultado principal (H1) \\ \hline
p75 de CLS & CrUX, \texttt{metrics\_summary} & Origen y mes & Variable de resultado secundaria \\ \hline
Rango de tráfico del origen & CrUX & Origen y mes & Covariable de emparejamiento y eje de heterogeneidad (H2) \\ \hline
País predominante & CrUX, tablas por país & Origen y mes & Covariable de emparejamiento \\ \hline
Cliente (móvil o escritorio) & HTTP Archive & Origen y mes & Estratificación: el panel se construye por separado para cada cliente \\ \hline
\end{tabular}
\end{table}

El tamaño del panel no crece de forma monótona (Figura \ref{fig:panel}): se identifican dos rampas de expansión fuerte (2018-2019 y 2022-2023) seguidas de un período de relativa estabilidad; desde mediados de 2023 la cantidad de orígenes cubiertos oscila estacionalmente entre aproximadamente 15,5 y 17 millones (mínimos en enero, máximos a mitad de año), sin una tendencia de crecimiento neta. Esta estabilidad reciente favorece construir un panel balanceado sobre la ventana 2023-2026, sin impedir extender el análisis a años anteriores siempre que el cambio de composición de la muestra de 2018-2019 se trate explícitamente como un quiebre estructural y no como parte de un panel homogéneo. En términos de costo computacional, extraer la dimensión de tecnología por origen y por mes insume aproximadamente 13 GB de escaneo en BigQuery por cada mes consultado, por lo que construir el panel completo de ocho años implica un escaneo del orden de 1,2 TB, alcanzable dentro de la cuota gratuita mensual de BigQuery distribuyendo la extracción en pocos meses o muestreando trimestralmente.

\begin{figure}[htbp]
\centering
\begin{tikzpicture}
\begin{axis}[
    width=0.92\textwidth, height=6cm,
    ylabel={Páginas raíz en el crawl (millones)},
    ylabel style={font=\small},
    ymin=0, ymax=19.5,
    xmin=0.4, xmax=14.6,
    xtick={1,...,14},
    xticklabels={01/20,07/20,01/21,07/21,01/22,07/22,01/23,07/23,
                 01/24,07/24,01/25,07/25,01/26,07/26},
    x tick label style={font=\scriptsize, rotate=45, anchor=north east},
    y tick label style={font=\scriptsize},
    grid=major, grid style={gray!20},
    axis lines=left,
    mark size=1.7pt,
]
\fill[itbablue!7] (axis cs:7,0) rectangle (axis cs:14.6,19.5);
\node[font=\scriptsize\itshape, itbablue, anchor=north east]
      at (axis cs:14.4,5.6) {ventana estable 2023--2026};
\addplot[color=itbablue, thick, mark=*] coordinates {
    (1,5.660) (2,6.975) (3,7.824) (4,8.217) (5,8.446) (6,10.787) (7,16.748)
    (8,17.174) (9,15.792) (10,16.745) (11,15.701) (12,16.273) (13,15.480) (14,16.268)
};
\end{axis}
\end{tikzpicture}
\caption{Páginas raíz únicas relevadas por HTTP Archive, en cortes semestrales de enero y julio. La rampa de 2022-2023 corresponde a la expansión del corpus; a partir de 2023 la serie se estabiliza y solo muestra la oscilación estacional entre el mínimo de enero y el máximo de julio.}
\label{fig:panel}
\end{figure}

\subsection{Chequeo exploratorio sobre dos cortes}
\label{sec:exploratorio}

Sobre los cortes de junio de 2025 y junio de 2026 se corrió un cálculo exploratorio, previo al trabajo de tesis, para ver si el dato se comporta como el diseño necesita. Para cada par de tecnologías se comparó la variación del p75 de LCP e INP entre junio de 2025 y junio de 2026 en los orígenes que migraron contra la variación en los orígenes que siguieron en la tecnología de origen. La diferencia entre ambas variaciones es un diferencias en diferencias de dos períodos, el caso más simple posible del estimador. Los resultados están en la Tabla \ref{tab:exploratorio}.

\begin{table}[H]
\centering
\small
\caption{Diferencias en diferencias exploratorio de dos períodos, junio de 2025 contra junio de 2026. Un valor negativo indica mejora del grupo migrado respecto de su control. Valores en milisegundos.}
\label{tab:exploratorio}
\begin{tabular}{|p{3.3cm}|r|r|c|r|r|}
\hline
\textbf{Par} & \textbf{Migrados} & \textbf{Control} & \textbf{LCP previo (mig./ctrl.)} & \textbf{DiD LCP} & \textbf{DiD INP} \\ \hline
jQuery $\to$ React        & 62.457 & 6.898.203 & 2.100 / 1.900 & $+41$  & $+41$ \\ \hline
Abandono de WordPress     & 65.407 &   627.691 & 2.400 / 2.100 & $-671$ & $+12$ \\ \hline
React $\to$ Next.js       &  9.059 &   751.129 & 2.200 / 1.900 & $-34$  & $0$   \\ \hline
\end{tabular}
\end{table}

Lo que aparece no es una confirmación de la hipótesis principal. En el par jQuery a React, el más central del trabajo, los sitios migrados mejoran su LCP en promedio 10\,ms mientras que los que se quedaron en jQuery mejoran 51\,ms: el grupo migrado queda 41\,ms por detrás de su control, y en INP queda 41\,ms por detrás también. Solo el abandono de WordPress muestra una mejora clara de LCP, y la migración de React a Next.js una mejora modesta. Si el cálculo terminara acá, la conclusión sería que migrar a React empeora la experiencia de usuario.

Hay dos razones para no leerlo así, y las dos son argumentos a favor del diseño que propone esta tesis. La primera es que los grupos no son comparables de entrada: en los tres pares, los sitios que migran arrancan entre 200 y 300\,ms peor que los que no migran, lo cual es esperable, porque un sitio lento tiene más motivos para migrar que uno rápido. Con esa diferencia de partida, buena parte de lo que se mide es reversión a la media y no efecto del tratamiento, y por eso el diseño de la tesis empareja el control por covariables previas al evento en vez de usar como control a todos los que no migraron. La segunda es que dos cortes anuales no permiten distinguir el momento del tratamiento: un sitio que migró en julio de 2025 y otro que migró en mayo de 2026 entran acá con el mismo peso, aunque uno tuvo once meses para que el cambio se notara y el otro uno. Esa confusión entre cohortes de adopción es exactamente el problema que resuelve el estimador de Callaway y Sant'Anna.

El ejercicio tampoco exige persistencia del cambio durante varios meses, con lo cual una parte de los eventos contados puede ser ruido de detección de Wappalyzer, y no permite verificar tendencias paralelas, porque con un solo período previo no hay tendencia que mirar. Su valor es otro: muestra que el cruce entre las dos fuentes produce grupos de tamaño suficiente, que la variable de resultado se mueve lo bastante como para que haya algo que estimar, y que la respuesta ingenua a la pregunta de investigación difiere de la que arroja un diseño con control adecuado. Ese contraste es, en sí mismo, parte de lo que la tesis se propone documentar.

\section{Hipótesis}

\textbf{Hipótesis principal (H1).} Las migraciones desde tecnologías frontend con mayor carga de JavaScript bloqueante hacia frameworks modernos (por ejemplo, de jQuery a React) producen una mejora estadísticamente significativa en el LCP y el INP de los sitios migrados, respecto de la trayectoria que hubieran seguido de no migrar, estimada a partir de un grupo de control emparejado. La evaluación se realiza mediante un diseño de estudio por eventos (event study) alrededor del mes de la migración, y se busca rechazar la hipótesis nula de ausencia de efecto. El chequeo exploratorio de la Sección \ref{sec:exploratorio} apunta en sentido contrario para el par jQuery a React, de modo que la hipótesis se formula sabiendo que el dato crudo no la respalda: parte del trabajo consiste en determinar si esa discrepancia se explica por la falta de emparejamiento y por la confusión entre cohortes de adopción, o si el efecto esperado sencillamente no existe.

\textbf{Hipótesis secundaria (H2).} El efecto de la migración no es homogéneo: es mayor en las migraciones que parten de tecnologías legacy (jQuery, plantillas de WordPress sin optimizar) que en las migraciones entre frameworks ya modernos (por ejemplo, de React a Next.js), y varía según el nivel de tráfico del sitio.

\textbf{Hipótesis secundaria (H3).} Parte del efecto que arrojaría un estimador clásico de efectos fijos bidireccionales (TWFE) sobre este panel está inflado por el sesgo de comparaciones entre grupos ya tratados que describe Goodman-Bacon \cite{goodmanbacon2021}; se espera que el estimador robusto de Callaway y Sant'Anna \cite{callawaysantanna2021} arroje un efecto de menor magnitud que el estimador clásico, aunque igualmente significativo.

Las variables operativas de estas hipótesis, tecnología detectada y percentil 75 de LCP, INP y CLS, fueron descriptas en la Sección \ref{sec:dataset}.

\section{Objetivos}

\subsection{Objetivo General}

Estimar el efecto causal de las migraciones de tecnología frontend sobre las métricas de experiencia de usuario real (Core Web Vitals), a partir de un panel longitudinal público construido con datos de HTTP Archive y CrUX, mediante técnicas de inferencia causal para datos de panel con tratamiento escalonado.

\subsection{Objetivos Específicos}

\begin{enumerate}
    \item \textbf{Construcción del panel:} construir un panel mensual de orígenes web con tecnología detectada y Core Web Vitals asociados, para una ventana temporal de al menos seis años.
    \item \textbf{Detección y validación de migraciones:} detectar eventos de migración tecnológica sobre un conjunto acotado de pares de tecnologías, y validar la precisión de esa detección.
    \item \textbf{Definición del grupo de control:} definir, para cada sitio migrado, un grupo de control emparejado de sitios comparables que no migraron.
    \item \textbf{Estimación del efecto causal:} estimar el efecto causal de la migración sobre LCP, INP y CLS mediante un diseño de diferencias en diferencias robusto a tratamiento escalonado.
    \item \textbf{Caracterización de heterogeneidad:} caracterizar cómo varía ese efecto según el tipo de migración y las características del sitio.
\end{enumerate}

\section{Solución propuesta a desarrollar en la tesis}
\label{sec:solucion}

La solución propuesta parte de una premisa simple: cada vez que un sitio cambia de tecnología frontend, ese cambio queda registrado mes a mes en HTTP Archive, y su efecto sobre la experiencia de los usuarios queda registrado mes a mes en CrUX. Cruzando ambas fuentes es posible tratar cada migración como un experimento natural y estimar su efecto con las mismas herramientas que se usan en economía aplicada para evaluar el impacto de una política pública sobre un conjunto de unidades que la reciben en distintos momentos.

La Figura \ref{fig:diseno} esquematiza el diseño para un sitio tratado y su control. La cantidad que se estima es la distancia, después del evento, entre la trayectoria observada del sitio que migró y la trayectoria contrafáctica que hubiera seguido de no migrar, aproximada por el grupo de control. El período previo al evento no aporta a la estimación, pero es donde se pone a prueba el supuesto que la sostiene: si tratados y controles no venían moviéndose en paralelo antes de la migración, no hay razón para creer que el control representa el contrafáctico después.

\begin{figure}[H]
\centering
\begin{tikzpicture}
\begin{axis}[
    width=0.92\textwidth, height=6.4cm,
    xmin=-4.5, xmax=4.6, ymin=2.15, ymax=3.65,
    xlabel={Meses respecto del evento de migración},
    ylabel={LCP p75 (menor es mejor)},
    xlabel style={font=\small}, ylabel style={font=\small},
    xtick={-4,-3,-2,-1,0,1,2,3},
    ytick=\empty,
    x tick label style={font=\scriptsize},
    axis lines=left,
    clip=false,
    legend style={at={(0.5,-0.32)}, anchor=north, legend columns=3,
                  font=\scriptsize, draw=gray!40},
]
\fill[gray!8] (axis cs:-4.5,2.15) rectangle (axis cs:0,3.65);
\draw[dashed, gray!70] (axis cs:0,2.15) -- (axis cs:0,3.55);
\node[font=\scriptsize\itshape, gray!65, anchor=north west] at (axis cs:-4.35,3.58)
      {ventana previa: chequeo de tendencias paralelas};

\addplot[color=itbablue, thick, mark=*, mark size=1.6pt] coordinates {
    (-4,3.30) (-3,3.26) (-2,3.21) (-1,3.16) (0,3.11)
    (1,2.72) (2,2.56) (3,2.50)
};
\addlegendentry{Sitio que migra}

\addplot[color=itbagris, thick, mark=square*, mark size=1.5pt] coordinates {
    (-4,3.00) (-3,2.96) (-2,2.91) (-1,2.86) (0,2.81)
    (1,2.77) (2,2.72) (3,2.68)
};
\addlegendentry{Control emparejado}

\addplot[color=itbablue, thick, dashed, forget plot] coordinates {
    (0,3.11) (1,3.07) (2,3.02) (3,2.98)
};
\addlegendimage{color=itbablue, thick, dashed}
\addlegendentry{Contrafáctico}

\draw[decorate, decoration={brace, amplitude=5pt}, itbablue, thick]
    (axis cs:3.18,2.98) -- (axis cs:3.18,2.50)
    node[midway, right=7pt, font=\scriptsize, itbablue] {ATT};
\end{axis}
\end{tikzpicture}
\caption{Esquema del diseño cuasi-experimental para un sitio tratado. Antes del evento, tratado y control evolucionan en paralelo; después, la diferencia entre la trayectoria observada del sitio que migró y su contrafáctico es el efecto que se busca estimar. Los valores son ilustrativos.}
\label{fig:diseno}
\end{figure}

\subsection{Metodología}

\textbf{(a) Construcción del panel.} Se extrae de HTTP Archive, mes a mes, la tecnología detectada para cada origen que además tenga métricas publicadas en CrUX ese mismo mes, separando cliente móvil y de escritorio. Se releva primero el rango de fechas disponible en ambas fuentes y la cobertura del cruce entre ellas, y se define la cohorte de orígenes elegibles de forma point-in-time, según su estado en el mes anterior a cualquier posible migración, para no introducir sesgo de supervivencia.

\textbf{(b) Detección y validación de eventos de migración.} Se define una migración como un origen que deja de tener detectada la tecnología A y pasa a tener detectada la tecnología B, y que sostiene ese cambio durante al menos tres meses consecutivos, para filtrar el ruido de detección. La dimensión de tecnología del panel se versiona como una dimensión de cambio lento (SCD tipo 2), donde cada fila registra un origen, una tecnología y un rango de vigencia; la fecha del evento de tratamiento es simplemente el inicio de vigencia del nuevo valor. Antes de elegir qué pares de tecnología estudiar, se analiza para cada tecnología su patrón de adopción y abandono a lo largo del tiempo, priorizando pares donde la detección es estable. Sobre una muestra aleatoria de migraciones detectadas se realiza una validación manual contra copias históricas del sitio (Wayback Machine) para estimar la precisión de la detección.

\textbf{(c) Diseño cuasi-experimental y grupo de control.} Para cada sitio migrado se construye un grupo de control de orígenes que no migraron en la misma ventana temporal, emparejados por covariables previas al evento: nivel de tráfico (rango de CrUX), país predominante y nivel de Core Web Vitals en los meses anteriores a la migración. Se verifica el supuesto de tendencias paralelas comparando la evolución de las métricas entre tratados y controles en el período previo al evento, y las covariables de emparejamiento se calculan siempre con información estrictamente anterior al mes de la migración de cada sitio. Este cuidado no es solo una precaución teórica: en un proyecto anterior de la maestría, sobre un problema de recomendación, una definición de corte temporal aparentemente razonable resultó filtrar información del futuro y produjo una mejora que solo existía dentro de la muestra de entrenamiento, sin sostenerse en una evaluación honesta. Ese antecedente es la razón concreta por la que acá el chequeo de tendencias paralelas se trata como un paso obligatorio y no como un formalismo.

\textbf{(d) Estimación del efecto causal.} Se estima un diseño de estudio por eventos alrededor del mes de la migración. Como el tratamiento ocurre en distintos meses para distintos sitios, se evita el estimador clásico de efectos fijos bidireccionales por el sesgo que señala Goodman-Bacon \cite{goodmanbacon2021}, y se utiliza en su lugar el estimador de Callaway y Sant'Anna \cite{callawaysantanna2021}, que estima el efecto por separado para cada cohorte de adopción y cada período, y lo agrega de forma consistente. Como verificación de robustez, se replica la estimación con el estimador de Sun y Abraham \cite{sunabraham2021}, que corrige el mismo sesgo por una vía distinta (una interacción saturada por cohorte y período en vez de una agregación por grupos); si ambos coinciden, refuerza la confianza en el resultado. Los errores estándar se calculan clusterizados por origen, y ante la comparación de varios pares de tecnología en simultáneo se aplica corrección por múltiples pruebas (Benjamini-Hochberg).

\textbf{(e) Heterogeneidad de efectos.} Sobre las migraciones con efecto estimado, se ajustan bosques aleatorios causales \cite{wagerathey2018} para estimar el efecto de tratamiento condicional a covariables como el tipo de migración, el nivel de tráfico previo y el país, sin imponer una forma funcional arbitraria a esa heterogeneidad, y se reporta la importancia relativa de cada covariable en explicar las diferencias de efecto.

\subsection{Técnicas / Herramientas}

El procesamiento del panel se realiza en BigQuery mediante SQL, dado el volumen de datos involucrado (varios terabytes de crawls mensuales). El análisis posterior se realiza en Python, con pandas y NumPy para la manipulación de datos. El ecosistema de Python para estimadores robustos a heterogeneidad de tratamiento bajo adopción escalonada es todavía más limitado que el de R, por lo que la estimación de Callaway y Sant'Anna se implementa directamente con el paquete \texttt{did} de R, invocado desde el pipeline de Python vía \texttt{rpy2}; el estimador de Sun y Abraham, usado como verificación de robustez, sí tiene una implementación madura en Python (\texttt{pyfixest}, función \texttt{sunab}) y se corre de forma nativa ahí. \texttt{statsmodels} se usa para los modelos de referencia (TWFE) y las pruebas de hipótesis. El emparejamiento del grupo de control se realiza con \texttt{scikit-learn}, usando un modelo de gradient boosting como estimador del puntaje de propensión y búsqueda de vecinos más cercanos sobre ese puntaje. Los bosques aleatorios causales se ajustan con la librería \texttt{econml} o, alternativamente, con \texttt{grf} de R vía \texttt{rpy2}. Para la reproducibilidad se usa control de versiones con Git y semillas aleatorias fijas en cada etapa estocástica del pipeline.

\section{Cronograma de actividades}

El cronograma arranca el 31 de agosto de 2026 y se extiende por cuarenta y ocho semanas. Las etapas están en correlato directo con los objetivos específicos: las semanas 1 a 16 cubren los objetivos 1 y 2 (construcción del panel y detección de migraciones), las semanas 17 a 28 el objetivo 3 (grupo de control y validación de supuestos), las semanas 29 a 38 el objetivo 4 (estimación y robustez) y las semanas 39 a 42 el objetivo 5 (heterogeneidad).

\vspace{0.3cm}
\begin{center}
\small
\begin{tabular}{|c|p{3.2cm}|p{6.4cm}|p{3.2cm}|}
\hline
\textbf{Semana} & \textbf{Fechas} & \textbf{Tareas} & \textbf{Entregable} \\ \hline
1--4 & Desde 31/08/2026 hasta 25/09/2026 & Revisión final de literatura, diseño técnico del panel y relevamiento de cobertura del cruce entre fuentes & En proceso \\ \hline
5--10 & Desde 28/09/2026 hasta 06/11/2026 & Construcción del panel HTTP Archive + CrUX (pipeline SQL, dimensión SCD-2 de tecnología) & Panel construido \\ \hline
11--16 & Desde 09/11/2026 hasta 18/12/2026 & Detección de eventos de migración y validación manual de una muestra contra Wayback Machine & Catálogo de eventos validado \\ \hline
17--24 & Desde 21/12/2026 hasta 12/02/2027 & Selección de pares tecnológicos y armado del grupo de control por puntaje de propensión & Diseño cuasi-experimental \\ \hline
25--28 & Desde 15/02/2027 hasta 12/03/2027 & Verificación del supuesto de tendencias paralelas y ajustes al emparejamiento & Validación de supuestos \\ \hline
29--34 & Desde 15/03/2027 hasta 23/04/2027 & Estimación DiD escalonado (Callaway y Sant'Anna) sobre jQuery$\to$React, con TWFE como referencia & Resultados H1 y H3 \\ \hline
35--38 & Desde 26/04/2027 hasta 21/05/2027 & Extensión a los pares restantes y réplica con el estimador de Sun y Abraham & Análisis de robustez \\ \hline
39--42 & Desde 24/05/2027 hasta 18/06/2027 & Heterogeneidad de efectos con bosques aleatorios causales & CATE por covariables \\ \hline
43--48 & Desde 21/06/2027 hasta 30/07/2027 & Redacción final de la tesis y preparación de la defensa & Tesis finalizada \\ \hline
\end{tabular}
\end{center}

\vspace{0.2cm}
\noindent El bloque de las semanas 17 a 24 abarca ocho semanas de calendario para una carga de trabajo de aproximadamente cuatro, porque cae sobre el receso de verano; el margen sobrante actúa además como amortiguador para el resto del plan. La estimación asume una dedicación sostenida de entre quince y veinte horas semanales.

\section{Limitaciones}

CrUX solo publica métricas de orígenes con tráfico de Chrome suficiente, lo que excluye a los sitios de baja audiencia y deja la muestra compuesta mayormente por sitios de mediana o alta audiencia. La detección de tecnología de Wappalyzer no es perfecta: algunas tecnologías muestran un patrón de adopción y abandono simultáneo mes a mes (lo que sugiere ruido de detección más que cambios reales), mientras que otras son mucho más estables; esto se releva empíricamente antes de elegir los pares de migración a estudiar. La migración de tecnología no es un evento aleatorio: los sitios que deciden migrar pueden diferir sistemáticamente de los que no lo hacen, lo que se aborda con matching sobre covariables previas al evento y con la verificación del supuesto de tendencias paralelas, aunque esto reduce el riesgo sin eliminarlo del todo. Por último, las reglas de detección de Wappalyzer cambiaron a lo largo de los años, por lo que un cambio de tecnología detectado puede en algunos casos reflejar un cambio en las reglas y no una migración real del sitio; esto se mitiga exigiendo persistencia del cambio durante varios meses consecutivos y validando manualmente una muestra de casos.

\bibliographystyle{unsrt}
\bibliography{bibliografia}

\end{document}
